\scanpage{69}

et une unique \deleted{\unsure{quadrique|quadratique}} \inserted{conique} circonscrite à $\Pi$ i.e. à $S$, savoir la cônique

\[C\ \text{d'équation}\qquad q_{0}(x)-1=0\]

Notons que $q_0$, considéré comme forme quadratique sur l'espace des translations de $E$, fait de $E$ un ``plan \uline{euclidien}'' -- la forme est en effet non dégénérée. Le carré de la longueur des cotés du polygône $\Pi$ est donné par

\[\begin{aligned}
q_0(s_1-s_0)&=q_0(\zeta-1,\zeta^{-1}-1)=(\zeta-1)(\zeta^{-1}-1)\\
&=2-\alpha
\end{aligned}\]

i.e.

\[q_0(a)=2-\alpha\]

On notera que c'est tjrs $\ne0$ i.e. \deleted{\unsure{$2-\alpha$}} $\alpha\ne2$, puisque $\zeta\ne1$. [On pourrait donc normaliser \deleted{\unsure{$q_0$|$q$}} $q$ \unsure{différemt|différemment}, en normalisant par $q_0(a)=1$. Dans le cas général, c'est cette normalisation-là qui sera la meilleure, comme nous verrons \dots.]

\deleted{\uline{Commentaires ``philosophiques''} On a trouvé, pour le gpe $D_{n}$ fixé, et des corps}

% archive sheet number: 68
